\oldpage[IV-2]{90}

\subsection{Dimensión de un preesquema algebraico.}

\begin{proposition}[5.2.1]
\label{IV.2.5.2.1-fr}
Sean $k$ un cuerpo, $X$ un preesquema irreducible localmente de tipo finito
sobre $k$, $\xi$ su punto genérico. Entonces $X$ es biequidimensional y se tiene
$\dim(X)=\operatorname{deg.tr}_k \boldsymbol{k}(\xi)$.
\end{proposition}

\begin{proof}
Todo anillo local $\mathcal O_x$ es un anillo local de una $k$-álgebra de tipo
finito, por tanto cociente de un anillo local regular $(\mathbf{0},17.3.9)$, y se sabe
$(\mathbf{0},16.5.12)$ que estos anillos son catenarios; por consiguiente $X$ es un
espacio catenario (5.1.5), y, como $X$ es irreducible, basta (5.1.1) demostrar
que, para \emph{todo punto cerrado} $x\in X$, se tiene
\begin{equation}
\dim(\mathcal O_x)=\operatorname{deg.tr}_k \boldsymbol{k}(\xi).
\tag{5.2.1.1}\label{IV.2.5.2.1.1-fr}
\end{equation}
Se puede evidentemente suponer que $X$ es reducido y afín, por tanto
\emph{íntegro}, con anillo $A$, un álgebra de tipo finito sobre $k$. Sea
$n=\operatorname{deg.tr}_k \boldsymbol{k}(\xi)$, siendo $\boldsymbol{k}(\xi)=K$ el cuerpo de fracciones
de $A$. Se sabe (Bourbaki, \emph{Alg. comm.}, cap.~V, \S~3,
n\textsuperscript{o}~1, th.~1) que existe una sub-$k$-álgebra
$B=k[t_1,\ldots,t_n]$ de $A$, donde los $t_i$ son algebraicamente independientes
sobre $k$, tal que $A$ es una $B$-álgebra \emph{finita}. Sea
$\mathfrak m=\mathfrak j_x$, que por hipótesis es un ideal maximal de $A$;
$\mathfrak n=B\cap\mathfrak m$ es, por tanto, un ideal maximal de $B$
(Bourbaki, \emph{Alg. comm.}, cap.~V, \S~2, n\textsuperscript{o}~1, prop.~1),
y $A_{\mathfrak m}$ es un anillo local de la $B_{\mathfrak n}$-álgebra
\emph{finita} $S^{-1}A$, donde $S=B\setmin\mathfrak n$; como $B_{\mathfrak n}$ es
integralmente cerrado y $S^{-1}A$ es íntegro, se tiene
$\dim(A_{\mathfrak m})=\dim(B_{\mathfrak n})$ $(\mathbf{0},16.1.6)$. Es posible, pues,
limitarse al caso en que $A=k[t_1,\ldots,t_n]$; se sabe entonces que
$k'=\boldsymbol{k}(x)$ es una extensión \emph{finita} de $k$ (\textbf I, 6.4.2); sea
$A'=k'[t_1,\ldots,t_n]$; teniendo en cuenta (\textbf I, 2.4.6 y 3.3.14), existe un
ideal maximal $\mathfrak m'$ de $A'$ situado sobre $\mathfrak m$ y tal que
$k'$ es el cuerpo residual de $A'_{\mathfrak m'}$. Como $A'$ es íntegro y
es una $A$-álgebra finita, el mismo razonamiento anterior muestra que
$\dim(A'_{\mathfrak m'})=\dim(A_{\mathfrak m})$, de modo que la cuestión se reduce
al caso en que $A/\mathfrak m=k$. Ahora bien, en este caso $\mathfrak m$ está
generado por los polinomios $t_i-a_i$ (donde $a_i\in k$, $1\leq i\leq n$, y
los $a_i$ son las imágenes canónicas de los $t_i$ en $A/\mathfrak m$).
Sustituyendo $t_i$ por $t_i-a_i$, se ve finalmente que la cuestión se reduce al
caso en que $\mathfrak m=\sum_{i=1}^nAt_i$. El completado del anillo local
$A_{\mathfrak m}$ es entonces el anillo de series formales
$k[[t_1,\ldots,t_n]]$; se sabe que tiene la misma dimensión que
$A_{\mathfrak m}$ $(\mathbf{0},16.2.4)$ y, por otra parte, la dimensión de
$k[[t_1,\ldots,t_n]]$ es igual a $n$ $(\mathbf{0},17.1.4)$; de ahí la conclusión.
\end{proof}

\begin{corollary}[5.2.2]
\label{IV.2.5.2.2-fr}
Para un preesquema $X$ localmente de tipo finito sobre un cuerpo $k$, $\dim(X)$
coincide con el número definido en \textup{(4.1.1)}.
\end{corollary}

\begin{proof}
En efecto, si los $X_\alpha$ son los subpreesquemas reducidos cuyos espacios
subyacentes son las componentes irreducibles de $X$, se tiene
$\dim(X)=\sup_\alpha(\dim(X_\alpha))$ $(\mathbf{0},14.1.2.1)$ y basta aplicar (5.2.1) a
los $X_\alpha$.
\end{proof}
